\section{Lean Formalization}
\label{app:lean}
\lstdefinelanguage{lean}{
  morekeywords={def,inductive,mutual,end,match,with,fun,forall,exists,theorem,structure,where,let,if,then,else,do,pure,abbrev,namespace,import,deriving,by,sorry},
  morecomment=[l]{--},
  sensitive=true
}
\lstset{breaklines=true}
\newcommand{\leanfile}[1]{\texttt{lean/Funnelling/#1}}
\newsavebox{\fitbox}
\newcommand{\fit}[1]{\sbox{\fitbox}{#1}\ifdim\wd\fitbox>\linewidth\resizebox{\linewidth}{!}{\usebox{\fitbox}}\else\usebox{\fitbox}\fi}
\lstdefinestyle{leanside}{language=lean,basicstyle=\scriptsize\ttfamily,breaklines=true,breakindent=1.5em,columns=fullflexible,aboveskip=0pt,belowskip=0pt}
\newenvironment{cmp}[1]{\par\medskip\noindent\begin{minipage}{\textwidth}\textbf{#1}\par\smallskip\noindent}{\end{minipage}\par}
\newenvironment{paperside}{\begin{minipage}[t]{0.45\textwidth}\vspace{0pt}\scriptsize}{\end{minipage}\hfill}
\newenvironment{leanside}{\begin{minipage}[t]{0.53\textwidth}\vspace{0pt}}{\end{minipage}}
\newcommand{\mapping}[1]{\par\smallskip\noindent{\scriptsize\textit{Mapping.} #1}\par\smallskip\hrule}

This appendix relates the formalism of Section~\ref{sec:core} and Section~\ref{sec:java} to a Lean~4 formalization in the folder \texttt{lean} of the paper repository.
The Lean formalization is set up, but the proofs of the soundness theorems are not done yet.

\subsection{Request}
This appendix answers the following request, reported verbatim, so that we can check later how much of it is addressed.
\begin{lstlisting}[language={},basicstyle=\small\ttfamily,breaklines=true,columns=fullflexible]

The Appendix, for each piece of formalism we report the following:

(A)Original formalism

(B)Corresponding Lean formalism

(C)English explanation of the differences and either
- why it ok for those differences to be there.
- or what the more appropriate form of the original formalism should be

(D)A clear report of the lean statements and an high level description of the proofs.

(E)A list of validation examples. Those are:
- programs that Do pass the lean type system AND the paper type system, or  that do not pass both the lean type system and the paper type system.
- programs that reduce in some steps equally in both the paper type system and the lean type system. Basically, this is 'testing': we are empirically checking that the two formalism are the same by checking that they behave the same on a range of examples including interesting corner cases.
\end{lstlisting}

\subsection{Status}
The file \texttt{lean/lean-toolchain} selects Lean 4.35.0: elan runs that version for every command in the folder \texttt{lean}.
The project depends on Mathlib v4.35.0, which brings Batteries and Aesop; \texttt{lake exe cache get} downloads Mathlib already compiled.
The current files import only the core library of Lean.
\texttt{lake build} in the folder \texttt{lean} checks all the files.
The formalization has two levels.
The \emph{high level}, \leanfile{High}, follows the metarules of the paper as closely as possible; Section~\ref{app:lean:sidebyside} compares it with the paper.
The \emph{low level}, \leanfile{Low}, is free: it can change in any way that makes the proofs easier, and several low levels can be tried, without changing the high level.
The low level starts as a copy of the high level, with its own datatypes and functions.
The files are:
\begin{itemize}
\item \leanfile{High/Syntax.lean}: types, funnels, expressions, methods, declarations, class declarations $K$.
\item \leanfile{High/Substitution.lean}: substitution of type variables and of variables, and the renaming used by $\mathsf{fresh}$.
\item \leanfile{High/Declarations.lean}: $\mathsf{declared}$, $\mathsf{scope}$, the pairing, $\mathrm{generic}$, $M\in C$, $\mathsf{typeOf}$, values, $\mathsf{fresh}$.
\item \leanfile{High/WellFormed.lean}: the well-formedness requirements A1-A11.
\item \leanfile{High/Typing.lean}: subtyping, kinding, checked casts, $\mathsf{overrideOk}$, $\mathsf{consistent}$, typing, method, class and program typing.
\item \leanfile{High/Reduction.lean}: evaluation contexts and the reduction.
\item \leanfile{Low/*.lean}: the same files for the low level, and \leanfile{Low/Soundness.lean}, progress and preservation of the low level.
\item \leanfile{Translation.lean}: the translation from the high level to the low level.
\item \leanfile{Adequacy.lean}: the lemmas that relate the two levels through the translation.
\item \leanfile{Soundness.lean}: the soundness statements of the paper, on the high level.
\item \leanfile{Examples.lean}: the validation examples, on the high level.
\end{itemize}
Lean accepts all the files.
The high-level \texttt{progress} and \texttt{preservation} are proved from the low-level ones through the adequacy lemmas.
The high-level \texttt{typeSoundness} is proved from \texttt{progress}, \texttt{preservation} and \texttt{stepWellFormed}.
The low-level \texttt{progress} and \texttt{preservation}, the six adequacy lemmas and \texttt{stepWellFormed} are stated with \texttt{sorry}.
Section~\ref{app:lean:examples} lists which validation examples are proved and which are stated with \texttt{sorry}.
The mini proof that checked the installation is not in the repository:
it proved $a+b=b+a$ with \texttt{omega}, $\mathsf{Even}\ 4$ for an inductive predicate $\mathsf{Even}$, $\neg\mathsf{Even}\ 1$ by \texttt{cases}, and evaluated an expression with \texttt{\#eval}.
\\[1ex]
The Lean code is written in ASCII:
\texttt{->} is $\rightarrow$, \texttt{/\char`\\} is $\wedge$, \texttt{\char`\\/} is $\vee$, \texttt{List.Mem x xs} is $x\in xs$, and \texttt{delta}, \texttt{gamma}, \texttt{sigma} are $\Delta$, $\Gamma$, $\sigma$.

\subsection{Differences Common to All the Pieces}
\label{app:lean:common}
The following differences hold for every piece below, and the pieces do not repeat them.
\begin{itemize}
\item \textbf{Names.}
Lean has one type \texttt{Name} (strings) for class names, method names, field names, variables and type variables.
The paper uses a different metavariable for each.
This is fine: no definition compares two names of different kinds.
\item \textbf{Grammar and well-formedness.}
In the paper, the metavariables range over the terms of the grammar, and the well-formedness requirements restrict the grammar (Section~\ref{sec:clarifications}).
In particular, a metavariable $v$ ranges over values and a metavariable $G$ ranges over ground types.
Lean defines the raw syntax as inductive types, and the restrictions as Boolean functions:
\texttt{Expr.isValue} for values, \texttt{Ty.isGround} for ground types, \texttt{wfDecls} and \texttt{wfExpr} for A1-A11.
Thus every Lean rule with a metavariable $v$ or $G$ of the paper has an explicit premise \texttt{isValue} or \texttt{allGround}.
The theorems assume \texttt{wfDecls} and \texttt{wfExpr} for the initial program.
The Lean reduction does not require the result to be well formed; the theorem \texttt{stepWellFormed} states that it is.
In the paper an instance of a reduction rule exists only when its result is an expression, so the paper needs no such lemma.
The two formulations agree when \texttt{stepWellFormed} holds.
\item \textbf{Sequences.}
The paper applies a judgement to sequences position by position, and requires the same length (Section~\ref{sec:declarations}).
Lean writes a list premise as a length equality and a premise for every pair in \texttt{List.zip}, or as a premise for every element of a list.
\item \textbf{Substitution and lengths.}
The paper defines $\tsubst{T}{\Delta}{\pl{T}}$ only when $\Delta$ and $\pl{T}$ have the same length.
Lean builds a substitution with \texttt{mkSubst xs ts}, which is \texttt{none} when the lengths differ; a premise \texttt{mkSubst xs ts = some sigma} requires the same length.
\item \textbf{Functions instead of relations.}
The paper defines $\mathsf{declared}$, $\mathsf{scope}$, $\mathrm{generic}$, the pairing and $\mathsf{typeOf}$ with \textbf{Define} boxes; those definitions have at most one result.
Lean defines them as functions, so the validation examples can compute them.
The inference rules (reduction, typing, kinding, subtyping, checked cast, method typing) are inductive relations in Lean, as in the paper.
\item \textbf{Equality of types.}
Lean derives no decidable equality for the nested inductive types; \texttt{Ty.beq} is a Boolean equality on types, used only by the pairing.
The proofs need the lemma \texttt{Ty.beq t t' = true <-> t = t'}.
\end{itemize}

\subsection{Side by Side: Paper and Lean}
\label{app:lean:sidebyside}
Each block below shows, on the left, a piece of the formalism of the paper, with the minimal text of the paper that the piece needs; on the right, the Lean definitions of the same concepts in the high level \texttt{Funnelling.High}; under them, the words that map one side to the other, and the differences beyond Section~\ref{app:lean:common}.

\subsubsection{Syntax}

\begin{cmp}{Types, funnels, bounds and type environments}
\begin{paperside}
\fit{$\begin{array}{rcl}
T & ::= & \gen{C}{\pl{T}} \mid X\\
F & ::= & G \mid \funnel{X}\\
G & ::= & \gen{C}{\pl{G}}\\
\Gamma & ::= & T_1\gap x_1,\ldots,T_n\gap x_n\\
\Delta & ::= & \ext{X_1}{T_1},\ldots,\ext{X_n}{T_n}
\end{array}$}\\[1ex]
A ground type $G$ is a type with no type variables.
For a type variable $X$, $\Delta(X)$ is the bound of $X$ in $\Delta$.
\end{paperside}
\begin{leanside}
\begin{lstlisting}[style=leanside]
inductive Ty where
  | cls (c : Name) (args : List Ty)
  | var (x : Name)
inductive Funnel where
  | ground (g : Ty)
  | funnel (x : Name)
abbrev Bounds := List (Prod Name Ty)
abbrev Ctx := List (Prod Name Ty)
def Ty.isGround : Ty -> Bool
  | .var _ => false
  | .cls _ ts => Ty.allGround ts
\end{lstlisting}
\end{leanside}
\mapping{$G$ is a \texttt{Ty} with \texttt{isGround}. A pair of $\Gamma$ is \texttt{(x, T)}: name first. $\Delta(X)$ is \texttt{delta.lookup X} and $\Gamma(x)$ is \texttt{gamma.lookup x}.}
\end{cmp}

\begin{cmp}{Expressions and methods}
\begin{paperside}
\fit{$\begin{array}{rcl}
e & ::= & x \mid \call{e}{\pl{T}}{m}{\pl{e}} \mid \explicitcall{e}{T}{\pl{T}}{m}{\pl{e}}\\
&& \mid e\fstop f \mid \obj{C}{\pl{T}}{\pl{e}}\\
&& \mid \innerobj{C}{\pl{F}}{C'}{\pl{T}'}{\pl{e}}{\pl{M}}\\
&& \mid \ifthen{e_0}{T}{x}{e_1}{e_2}\\
M & ::= & \method{\Delta}{T_0}{m}{T_1\gap x_1,\ldots,T_n\gap x_n}{\ret{e}}
\end{array}$}
\end{paperside}
\begin{leanside}
\begin{lstlisting}[style=leanside]
mutual
inductive Expr where
  | var (x : Name)
  | call (recv : Expr) (targs : List Ty) (m : Name) (args : List Expr)
  | superCall (recv : Expr) (start : Ty) (targs : List Ty) (m : Name) (args : List Expr)
  | field (recv : Expr) (f : Name)
  | new (c : Name) (targs : List Ty) (args : List Expr)
  | newLocal (c : Name) (fs : List Funnel) (sup : Name) (supArgs : List Ty)
      (args : List Expr) (ms : List Method)
  | ifInst (e0 : Expr) (t : Ty) (x : Name) (e1 e2 : Expr)
inductive Method where
  | mk (bounds : Bounds) (ret : Ty) (m : Name) (params : List (Prod Ty Name)) (body : Expr)
end
\end{lstlisting}
\end{leanside}
\mapping{The super call $\explicitcall{e}{T}{\pl{T}}{m}{\pl{e}}$ is \texttt{superCall}; the local class is \texttt{newLocal}; the type test is \texttt{ifInst}. A parameter $T_i\gap x_i$ is the pair \texttt{(T, x)}.}
\end{cmp}

\begin{cmp}{Declarations and programs}
\begin{paperside}
\fit{$\begin{array}{rcl}
\mathsf{Dec} & ::= & \subclass{C}{\Delta}{C'}{\pl{T}}{\pl{M}}{\pl{S}}\\
S & ::= & T\gap f\semi\\
K & ::= & \mathsf{Dec} \mid \innerobj{C}{\pl{F}}{C'}{\pl{T}}{\pl{e}}{\pl{M}}
\end{array}$}\\[1ex]
A program is a class table \decls and a main expression $e$.
\end{paperside}
\begin{leanside}
\begin{lstlisting}[style=leanside]
structure Dec where
  name : Name
  bounds : Bounds
  sup : Name
  supArgs : List Ty
  methods : List Method
  fields : List (Prod Ty Name)
inductive K where
  | top (dec : Dec)
  | loc (c : Name) (fs : List Funnel) (sup : Name) (supArgs : List Ty)
      (args : List Expr) (ms : List Method)
abbrev Decls := List Dec
\end{lstlisting}
\end{leanside}
\mapping{The slots $\pl{S}$ are \texttt{fields}, pairs \texttt{(T, f)}. A program is a pair \texttt{(d : Decls)} and \texttt{(e : Expr)}, passed separately.}
\end{cmp}

\begin{cmp}{Values}
\begin{paperside}
\fit{$v ::= \obj{C}{\pl{G}}{\pl{v}} \mid \innerobj{C}{\pl{G}}{C'}{\pl{G}'}{\pl{v}}{\pl{M}}$}
\end{paperside}
\begin{leanside}
\begin{lstlisting}[style=leanside]
def Expr.isValue : Expr -> Bool
  | .new _ gs vs => Ty.allGround gs && Expr.allValues vs
  | .newLocal _ fs _ gs vs _ =>
    fs.all Funnel.isGround && Ty.allGround gs && Expr.allValues vs
  | _ => false
\end{lstlisting}
\end{leanside}
\mapping{The funnels $\pl{G}$ of a value of a local class are ground funnels: \texttt{fs.all Funnel.isGround}.}
\end{cmp}

\begin{cmp}{Evaluation contexts}
\begin{paperside}
\fit{$\begin{array}{rcl}
\mathcal{E} & ::= & \square \mid \call{\mathcal{E}}{\pl{G}}{m}{\pl{e}} \mid \call{v}{\pl{G}}{m}{\pl{v},\mathcal{E},\pl{e}}\\
&& \mid \explicitcall{\mathcal{E}}{G}{\pl{G}}{m}{\pl{e}}\\
&& \mid \explicitcall{v}{G}{\pl{G}}{m}{\pl{v},\mathcal{E},\pl{e}}\\
&& \mid \mathcal{E}\fstop f \mid \obj{C}{\pl{G}}{\pl{v},\mathcal{E},\pl{e}}\\
&& \mid \innerobj{C}{\pl{G}}{C'}{\pl{G}'}{\pl{v},\mathcal{E},\pl{e}}{\pl{M}}\\
&& \mid \ifthen{\mathcal{E}}{G}{x}{e_1}{e_2}
\end{array}$}
\end{paperside}
\begin{leanside}
\begin{lstlisting}[style=leanside]
inductive ECtx where
  | hole
  | callRecv (ctx : ECtx) (gs : List Ty) (m : Name) (es : List Expr)
  | callArg (v : Expr) (gs : List Ty) (m : Name) (vs : List Expr) (ctx : ECtx) (es : List Expr)
  | superRecv (ctx : ECtx) (g : Ty) (gs : List Ty) (m : Name) (es : List Expr)
  | superArg (v : Expr) (g : Ty) (gs : List Ty) (m : Name) (vs : List Expr) (ctx : ECtx) (es : List Expr)
  | field (ctx : ECtx) (f : Name)
  | newArg (c : Name) (gs : List Ty) (vs : List Expr) (ctx : ECtx) (es : List Expr)
  | localArg (c : Name) (fs : List Funnel) (c' : Name) (gs : List Ty) (vs : List Expr) (ctx : ECtx)
      (es : List Expr) (ms : List Method)
  | ifCond (ctx : ECtx) (g : Ty) (x : Name) (e1 e2 : Expr)
\end{lstlisting}
\end{leanside}
\mapping{$\mathcal{E}[e]$ is \texttt{ECtx.plug E e}. \texttt{ECtx.ok E} checks that the $v$ are values and the $G$ are ground.}
\end{cmp}

\subsubsection{Well-Formedness}

\begin{cmp}{A1, A2, A3}
\begin{paperside}
(A1) In \decls, class names are unique across top-level and local classes.\\
(A2) In a class, all the field names, including the field names of the superclasses, are disjoint.\\
(A3) In a class, top-level or local, all the method names are disjoint.
\end{paperside}
\begin{leanside}
\begin{lstlisting}[style=leanside]
def a1 (d : Decls) : Bool := nodup ((declaredKs d).map K.name)
def a2 (d : Decls) : Bool :=
  d.all fun dec => nodup ((superChain d (d.length + 1) dec.name).flatMap
    fun c => c.fields.map Prod.snd)
def a3 (ks : List K) : Bool := ks.all fun k => nodup (k.methods.map Method.name)
\end{lstlisting}
\end{leanside}
\mapping{\texttt{superChain d n c} is the chain of superclasses of $c$, cut after $n$ classes; an acyclic chain is never longer than $|\decls|$, and A9 rejects a cyclic one. A3 is checked on \texttt{declaredKs d} and on the local classes of the main expression.}
\end{cmp}

\begin{cmp}{A4, A5, A6}
\begin{paperside}
(A4) In a method, all the parameter names $x_1,\ldots,x_n$ are disjoint.\\
(A5) In a $\Delta$, all the $X$ in $\mathrm{dom}(\Delta)$ are disjoint.\\
(A6) In a list of funnels $\pl{F}$, all the $X$ in $\mathrm{dom}(\pl{F})$ are disjoint.
\end{paperside}
\begin{leanside}
\begin{lstlisting}[style=leanside]
def a4 (ms : List Method) : Bool := ms.all fun m => nodup m.params
def a5 (d : Decls) (ms : List Method) : Bool :=
  d.all (fun dec => nodup (dom dec.bounds)) && ms.all fun m => nodup (dom m.bounds)
def a6 (ks : List K) : Bool := ks.all fun k => nodup (domF k.funnels)
\end{lstlisting}
\end{leanside}
\mapping{Each check runs on all the methods and local classes of \decls, and of the main expression.}
\end{cmp}

\begin{cmp}{A7}
\begin{paperside}
(A7) There is no shadowing of $X$ or $x$, except for \Q@this@.\\[1ex]
A method body of a local class also binds \Q@this@, thus shadowing the \Q@this@ of the enclosing method; this is the only form of shadowing.
\end{paperside}
\begin{leanside}
\begin{lstlisting}[style=leanside]
def Expr.noShadow (tvs vs : List Name) : Expr -> Bool
  ...
  | .newLocal _ _ _ _ es ms =>
    Expr.noShadowList tvs vs es &&
    Method.noShadowList tvs (vs.filter fun x => x != thisName) ms
  | .ifInst e0 _ x e1 e2 =>
    e0.noShadow tvs vs && !(vs.contains x) &&
    e1.noShadow tvs (x :: vs) && e2.noShadow tvs vs
def Method.noShadow (tvs vs : List Name) : Method -> Bool
  | .mk delta _ _ ps e =>
    (dom delta).all (fun x => !(tvs.contains x)) &&
    (ps.map Prod.snd).all (fun x => !(vs.contains x)) &&
    e.noShadow (tvs ++ dom delta) (thisName :: vs ++ ps.map Prod.snd)
def a7Decls (d : Decls) : Bool :=
  d.all fun dec => Method.noShadowList (dom dec.bounds) [] dec.methods
\end{lstlisting}
\end{leanside}
\mapping{\texttt{tvs} and \texttt{vs} are the type variables and the variables in scope. The cases written \texttt{...} go into the subexpressions with the same \texttt{tvs} and \texttt{vs}. Funnels bind nothing: $\funnel{X}$ refers to the $X$ in scope.}
\end{cmp}

\begin{cmp}{A8, A9}
\begin{paperside}
(A8) No parameter and no variable of an \Q@if@ is \Q@this@.\\
(A9) \decls contains $\subclass{\Object}{}{\Object}{}{}{}$, and the chain of superclasses of every other top-level class reaches \Q@Object@.
\end{paperside}
\begin{leanside}
\begin{lstlisting}[style=leanside]
def a8 (ms : List Method) (ifVars : List Name) : Bool :=
  ms.all (fun m => !(m.params.contains thisName)) && !(ifVars.contains thisName)
def a9 (d : Decls) : Bool :=
  (d.any fun dec => dec.name == objectName && dec.bounds.isEmpty &&
    dec.sup == objectName && dec.supArgs.isEmpty && dec.methods.isEmpty &&
    dec.fields.isEmpty) &&
  d.all fun dec => reachesObject d dec.name
\end{lstlisting}
\end{leanside}
\mapping{\texttt{ifVars} lists the variables of all the \Q@if@ in the program.}
\end{cmp}

\begin{cmp}{A10, A11}
\begin{paperside}
(A10) In \decls, the superclass of every class, top-level or local, is a top-level class.\\
(A11) In \decls, every funnel is of the form $\funnel{X}$.
\end{paperside}
\begin{leanside}
\begin{lstlisting}[style=leanside]
def a10 (d : Decls) : Bool := (declaredKs d).all fun k => (findDec d k.sup).isSome
def a11 (d : Decls) : Bool :=
  (declaredKs d).all fun k => k.funnels.all fun
    | .funnel _ => true
    | .ground _ => false
def groundFunnels (ks : List K) : Bool :=
  ks.all fun k => k.funnels.all fun
    | .funnel _ => true
    | .ground g => g.isGround
\end{lstlisting}
\end{leanside}
\mapping{\texttt{groundFunnels} is the grammar $G::=\gen{C}{\pl{G}}$ of a ground funnel, not a requirement of the paper.}
\end{cmp}

\begin{cmp}{All the requirements}
\begin{paperside}
We require the well-formedness conditions A1-A11.\\[1ex]
The well-formedness requirements A1-A11 are part of the grammar (Section~\ref{sec:clarifications}).
\end{paperside}
\begin{leanside}
\begin{lstlisting}[style=leanside]
def wfDecls (d : Decls) : Bool :=
  a1 d && a2 d && a3 (declaredKs d) && a4 d.methods && a5 d d.methods &&
  a6 (declaredKs d) && a7Decls d &&
  a8 d.methods (d.flatMap fun dec => Method.ifVarsList dec.methods) &&
  a9 d && a10 d && a11 d && groundFunnels (declaredKs d)
def wfExpr (e : Expr) : Bool :=
  a3 e.locals && a4 e.methods && a5 [] e.methods && a6 e.locals &&
  e.noShadow [] [] && a8 e.methods e.ifVars && groundFunnels e.locals
\end{lstlisting}
\end{leanside}
\mapping{\texttt{wfDecls} checks \decls; \texttt{wfExpr} checks the main expression, which A3-A8 also restrict. The theorems assume both (Section~\ref{app:lean:common}).}
\end{cmp}

\subsubsection{Declarations}

\begin{cmp}{$\declared{K}$}
\begin{paperside}
\fit{$\begin{array}{@{}l@{}}
\declared{\mathsf{Dec}}\textrm{ iff }\mathsf{Dec}\in\decls\\[1ex]
\declared{\innerobj{C_0}{\pl{F}}{C'}{\pl{T}}{\pl{e}}{\pl{M}}}\\
\quad\textrm{where }\declared{K}\\
\quad\phantom{\textrm{where }}\method{\_}{\_}{m}{\_}{\ret{e}}\in K\\
\quad\phantom{\textrm{where }}\innerobj{C_0}{\pl{F}}{C'}{\pl{T}}{\pl{e}}{\pl{M}}\textrm{ occurs in }e
\end{array}$}\\[1ex]
An expression $e'$ occurs in $e$ when $e=\mathcal{C}[e']$, for an expression context $\mathcal{C}$.
Expression contexts $\mathcal{C}$ do not enter the methods of a local class.
\end{paperside}
\begin{leanside}
\begin{lstlisting}[style=leanside]
def Expr.locals : Expr -> List K
  | .var _ => []
  | .call e _ _ es => e.locals ++ Expr.localsList es
  | .superCall e _ _ _ es => e.locals ++ Expr.localsList es
  | .field e _ => e.locals
  | .new _ _ es => Expr.localsList es
  | .newLocal c fs c' ts es ms =>
    K.loc c fs c' ts es ms :: (Expr.localsList es ++ Method.localsList ms)
  | .ifInst e0 _ _ e1 e2 => e0.locals ++ e1.locals ++ e2.locals
def declaredKs (d : Decls) : List K :=
  d.flatMap fun dec => K.top dec :: Method.localsList dec.methods
def Declared (d : Decls) (k : K) : Prop := List.Mem k (declaredKs d)
\end{lstlisting}
\end{leanside}
\mapping{\texttt{Expr.locals} also enters the methods of the local classes it finds, so \texttt{declaredKs} lists in one pass all the local classes that the paper reaches through a chain of $\mathsf{declared}$.}
\end{cmp}

\begin{cmp}{$\mathsf{scope}(\gen{C}{\pl{T}})$}
\begin{paperside}
\fit{$\begin{array}{@{}l@{}}
\mathsf{scope}(\gen{C}{\pl{T}}) = \Delta;\pl{T}\\
\quad\textrm{where }\declared{\subclass{C}{\Delta}{\_}{\_}{\_}{\_}}\\
\mathsf{scope}(\gen{C_0}{\pl{T}}) = \pairing{\Delta,\Delta'}{\pl{F}}{\pl{T}}\\
\quad\textrm{where }\method{\Delta'}{\_}{m}{\_}{\ret{e}}\in C\\
\quad\phantom{\textrm{where }}\innerobj{C_0}{\pl{F}}{\_}{\_}{\_}{\_}\textrm{ occurs in }e\\
\quad\phantom{\textrm{where }}\mathsf{scope}(\generic{C})=\Delta;\_
\end{array}$}
\end{paperside}
\begin{leanside}
\begin{lstlisting}[style=leanside]
def Expr.scopeEntries (env : Bounds) : Expr -> List ScopeEntry
  ...
  | .newLocal c fs _ _ es ms =>
    (c, (env, fs)) :: (Expr.scopeEntriesList env es ++
      match pairing env fs (fs.map Funnel.abs) with
      | some (inner, _) => Method.scopeEntriesList inner ms
      | none => [])
def Method.scopeEntries (env : Bounds) : Method -> List ScopeEntry
  | .mk delta _ _ _ e => e.scopeEntries (env ++ delta)
def scope (d : Decls) (c : Name) (ts : List Ty) : Option (Prod Bounds (List Ty)) :=
  match findDec d c with
  | some dec => some (dec.bounds, ts)
  | none =>
    match (d.flatMap fun dec => Method.scopeEntriesList dec.bounds dec.methods).lookup c with
    | some (env, fs) => pairing env fs ts
    | none => none
\end{lstlisting}
\end{leanside}
\mapping{\texttt{env} is $\Delta,\Delta'$: the bounds in scope where the local class is written. In the methods of a local class, \texttt{inner} is the bounds of $\mathsf{scope}(\generic{C})$. By A1, \texttt{lookup} finds the only entry of $C_0$.}
\end{cmp}

\begin{cmp}{Pairing}
\begin{paperside}
\fit{$\begin{array}{@{}l@{}}
\pairing{\Delta}{\emptyset}{\emptyset} = \emptyset;\emptyset\\
\pairing{\Delta}{\funnel{X},\pl{F}}{T,\pl{T}} = (\ext{X}{\Delta(X)};T)\cdot(\pairing{\Delta}{\pl{F}}{\pl{T}})\\
\pairing{\Delta}{G,\pl{F}}{G,\pl{T}} = \pairing{\Delta}{\pl{F}}{\pl{T}}\\[1ex]
(\Delta;\pl{T})\cdot(\Delta';\pl{T}') = \Delta,\Delta';\pl{T},\pl{T}'
\end{array}$}
\end{paperside}
\begin{leanside}
\begin{lstlisting}[style=leanside]
def pairing (delta : Bounds) : List Funnel -> List Ty -> Option (Prod Bounds (List Ty))
  | [], [] => some ([], [])
  | .funnel x :: fs, t :: ts => do
    let b <- delta.lookup x
    let (delta', ts') <- pairing delta fs ts
    pure ((x, b) :: delta', t :: ts')
  | .ground g :: fs, t :: ts => if g.beq t then pairing delta fs ts else none
  | _, _ => none
\end{lstlisting}
\end{leanside}
\mapping{An undefined pairing is \texttt{none}. The operator $\cdot$ is the cons of the recursive case.}
\end{cmp}

\begin{cmp}{$M\in K$, $M\in C$, $\generic{C}$, $|F|$}
\begin{paperside}
\fit{$\begin{array}{@{}l@{}}
M\in K\textrm{ iff }M\in\pl{M}\\
\quad\textrm{where }K\in\{\subclass{\_}{\_}{\_}{\_}{\pl{M}}{\_},\innerobj{\_}{\_}{\_}{\_}{\_}{\pl{M}}\}\\
M\in C\textrm{ iff }M\in K\\
\quad\textrm{where }\declared{K}\\
\quad\phantom{\textrm{where }}K\in\{\subclass{C}{\_}{\_}{\_}{\_}{\_},\innerobj{C}{\_}{\_}{\_}{\_}{\_}\}\\[1ex]
\generic{C} = \gen{C}{\mathrm{dom}(\Delta)}\\
\quad\textrm{where }\declared{\subclass{C}{\Delta}{\_}{\_}{\_}{\_}}\\
\generic{C} = \gen{C}{\mathrm{dom}(\pl{F})}\\
\quad\textrm{where }\declared{\innerobj{C}{\pl{F}}{\_}{\_}{\_}{\_}}\\[1ex]
|\funnel{X}| = X\qquad|G| = G
\end{array}$}
\end{paperside}
\begin{leanside}
\begin{lstlisting}[style=leanside]
def K.methods : K -> List Method
  | .top dec => dec.methods
  | .loc _ _ _ _ _ ms => ms
def MethodIn (d : Decls) (mt : Method) (c : Name) : Prop :=
  exists k, Declared d k /\ k.name = c /\ List.Mem mt k.methods
def generic (d : Decls) (c : Name) : Option (List Name) :=
  match findK d c with
  | some (.top dec) => some (dom dec.bounds)
  | some (.loc _ fs _ _ _ _) => some (domF fs)
  | none => none
def Funnel.abs : Funnel -> Ty
  | .ground g => g
  | .funnel x => .var x
\end{lstlisting}
\end{leanside}
\mapping{$M\in K$ is \texttt{List.Mem M k.methods}; $M\in C$ is \texttt{MethodIn d M C}. \texttt{generic} returns the type variables $\pl{X}$ of $\gen{C}{\pl{X}}$. $|F|$ is \texttt{Funnel.abs}.}
\end{cmp}

\begin{cmp}{$\mathsf{typeOf}(v)$}
\begin{paperside}
\fit{$\begin{array}{@{}l@{}}
\mathsf{typeOf}(v)=\gen{C}{\pl{G}}\\
\quad\textrm{where }v\in\{\obj{C}{\pl{G}}{\_},\innerobj{C}{\pl{G}}{\_}{\_}{\_}{\_}\}
\end{array}$}
\end{paperside}
\begin{leanside}
\begin{lstlisting}[style=leanside]
def typeOf : Expr -> Option Ty
  | .new c gs _ => some (.cls c gs)
  | .newLocal c fs _ _ _ _ => some (.cls c (fs.map Funnel.abs))
  | _ => none
\end{lstlisting}
\end{leanside}
\mapping{For a value of a local class, the ground funnels $\pl{G}$ become types through \texttt{Funnel.abs}.}
\end{cmp}

\subsubsection{Substitution and Renaming}

\begin{cmp}{Substitution of type variables}
\begin{paperside}
\fit{$\begin{array}{@{}lcl@{}}
\tsubst{\funnel{X'}}{\pl{X}}{\pl{T}} & = & \funnel{X'} \quad\textrm{if }X'\not\in\pl{X}\\
\tsubst{\funnel{X_i}}{X_1,\ldots,X_n}{T_1,\ldots,T_n} & = & T_i\\
\tsubst{G}{\pl{X}}{\pl{T}} & = & G\\
\tsubst{X'}{\pl{X}}{\pl{T}} & = & X' \quad\textrm{if }X'\not\in\pl{X}\\
\tsubst{X_i}{X_1,\ldots,X_n}{T_1,\ldots,T_n} & = & T_i
\end{array}$}\\[1ex]
The other cases go into every subterm, including the methods of a local class.
In the domain of a substitution, bounds $\Delta$ stand for $\mathrm{dom}(\Delta)$, and funnels $\pl{F}$ stand for their funnelled type variables; the ground funnels are ignored.
Substitution of type variables into expressions happens only during reduction, where the substituted types are ground.
\end{paperside}
\begin{leanside}
\begin{lstlisting}[style=leanside]
def mkSubst (xs : List Name) (ts : List Ty) : Option TSubst :=
  if xs.length = ts.length then some (xs.zip ts) else none
def mkSubstF (fs : List Funnel) (ts : List Ty) : Option TSubst :=
  if fs.length = ts.length then
    some ((fs.zip ts).filterMap fun
      | (.funnel x, t) => some (x, t)
      | (.ground _, _) => none)
  else none
def Ty.subst (sigma : TSubst) : Ty -> Ty
  | .var x => match sigma.lookup x with
    | some t => t
    | none => .var x
  | .cls c ts => .cls c (Ty.substList sigma ts)
def Funnel.subst (sigma : TSubst) : Funnel -> Funnel
  | .ground g => .ground (g.subst sigma)
  | .funnel x => match sigma.lookup x with
    | some t => .ground t
    | none => .funnel x
\end{lstlisting}
\end{leanside}
\mapping{$[\Delta=\pl{T}]$ is \texttt{mkSubst (dom delta) ts}, and $[\pl{F}=\pl{T}]$ is \texttt{mkSubstF fs ts}; both are \texttt{none} when the lengths differ. \texttt{Expr.tsubst} and \texttt{Method.tsubst} go into every subterm. \texttt{Funnel.subst} does not check that the result is ground; \texttt{stepWellFormed} states it.}
\end{cmp}

\begin{cmp}{Substitution of variables}
\begin{paperside}
\fit{$\begin{array}{@{}l@{}}
\tsubst{x_i}{x_1,\ldots,x_n}{v_1,\ldots,v_n} = v_i\\[1ex]
\tsubst{(\innerobj{C}{\pl{F}}{C'}{\pl{T}}{\pl{e}}{\pl{M}})}{\pl{x}}{\pl{v}}\\
\quad = \innerobj{C}{\pl{F}}{C'}{\pl{T}}{\tsubst{\pl{e}}{\pl{x}}{\pl{v}}}{\tsubst{\pl{M}}{\pl{x}}{\pl{v}}}
\quad\textrm{if }\this\not\in\pl{x}\\[1ex]
\tsubst{(\innerobj{C}{\pl{F}}{C'}{\pl{T}}{\pl{e}}{\pl{M}})}{\pl{x},\this,\pl{x}'}{\pl{v},v,\pl{v}'}\\
\quad = \innerobj{C}{\pl{F}}{C'}{\pl{T}}{\tsubst{\pl{e}}{\pl{x},\this,\pl{x}'}{\pl{v},v,\pl{v}'}}{\tsubst{\pl{M}}{\pl{x},\pl{x}'}{\pl{v},\pl{v}'}}
\end{array}$}\\[1ex]
The other cases go into every subterm.
\end{paperside}
\begin{leanside}
\begin{lstlisting}[style=leanside]
def Expr.vsubst (sigma : VSubst) : Expr -> Expr
  | .var x => match sigma.lookup x with
    | some v => v
    | none => .var x
  ...
  | .newLocal c fs c' ts es ms =>
    .newLocal c fs c' ts (Expr.vsubstList sigma es)
      (Method.vsubstList (sigma.filter fun p => p.1 != thisName) ms)
def Method.vsubst (sigma : VSubst) : Method -> Method
  | .mk delta t m ps e => .mk delta t m ps (e.vsubst sigma)
\end{lstlisting}
\end{leanside}
\mapping{The two cases of the local class are one case: removing \Q@this@ from the substitution changes nothing when \Q@this@ is not in it.}
\end{cmp}

\begin{cmp}{$\mathsf{fresh}(e,e')$ and $\mathsf{fresh}(M,M')$}
\begin{paperside}
The names bound in an expression $e$ are the variable of every \Q@if@ occurring in $e$, and the names bound in every method of a local class occurring in $e$.
The names bound in a method $M$ are the type variables of its signature, its parameters, and the names bound in its body.
$\mathsf{fresh}(e,e')$ holds when $e'$ is $e$ where every name bound in $e$ is renamed into a fresh name, that is a name occurring nowhere in the program; distinct binders get distinct fresh names.
$\mathsf{fresh}(M,M')$ holds when $M'$ is $M$ renamed in the same way.
The variable \Q@this@ is never renamed.
\end{paperside}
\begin{leanside}
\begin{lstlisting}[style=leanside]
def freshNames (avoid ns : List Name) : Bool :=
  decide ns.Nodup && ns.all fun n => !(avoid.contains n)
def Fresh (avoid : List Name) (e e' : Expr) : Prop :=
  exists ns, Expr.freshen [] [] ns e = some (e', []) /\
    freshNames avoid ns = true
def FreshM (avoid : List Name) (mt mt' : Method) : Prop :=
  exists ns, Method.freshen [] [] ns mt = some (mt', []) /\
    freshNames avoid ns = true
def avoidNames (d : Decls) (e : Expr) : List Name := d.names ++ e.names
\end{lstlisting}
\end{leanside}
\mapping{\texttt{Expr.freshen rv rt ns e} walks $e$, takes the next name of \texttt{ns} at every binder, and renames the bound occurrences: \texttt{rv} renames variables, \texttt{rt} type variables. The rules pass \texttt{avoidNames d redex}: the names of \decls and of the redex, fewer than all the names of the program. This is enough for A7, since the hole of an evaluation context is under no binder.}
\end{cmp}

\subsubsection{Reduction}

\begin{cmp}{\textsc{simpleFieldFound}}
\begin{paperside}
\fit{$\inferrule{
    \declared{\subclass{C}{\_}{\_}{\_}{\_}{T_1\gap f_1\semi\ldots T_n\gap f_n\semi}}\\
    f_i\in f_1,\ldots,f_n
}{
    \step{\obj{C}{\pl{G}}{v_1,\ldots,v_n,\pl{v}}\fstop f_i}{v_i}
}$}
\end{paperside}
\begin{leanside}
\begin{lstlisting}[style=leanside]
| simpleFieldFound {c gs vs1 vs f v} {dec : Dec} :
    List.Mem dec d -> dec.name = c ->
    vs1.length = dec.fields.length ->
    (Expr.new c gs (vs1 ++ vs)).isValue = true ->
    ((fieldNames dec).zip vs1).lookup f = some v ->
    Step d (.field (.new c gs (vs1 ++ vs)) f) v
\end{lstlisting}
\end{leanside}
\mapping{$v_1,\ldots,v_n$ is \texttt{vs1}, as long as the fields; $f_i$ and $v_i$ are found by \texttt{lookup} in the pairs of field names and values, which A2 makes unique.}
\end{cmp}

\begin{cmp}{\textsc{simpleFieldSuper}}
\begin{paperside}
\fit{$\inferrule{
    \declared{\subclass{C}{\Delta}{C'}{\pl{T}}{\_}{T_1\gap f_1\semi\ldots T_n\gap f_n\semi}}\\
    f\not\in f_1,\ldots,f_n\\
    C\neq\Object
}{
    \step{\obj{C}{\pl{G}}{v_1,\ldots,v_n,\pl{v}}\fstop f}{\obj{C'}{\tsubst{\pl{T}}{\Delta}{\pl{G}}}{\pl{v}}\fstop f}
}$}
\end{paperside}
\begin{leanside}
\begin{lstlisting}[style=leanside]
| simpleFieldSuper {c gs vs1 vs f sigma} {dec : Dec} :
    List.Mem dec d -> dec.name = c ->
    vs1.length = dec.fields.length ->
    (Expr.new c gs (vs1 ++ vs)).isValue = true ->
    (fieldNames dec).contains f = false ->
    Not (c = objectName) ->
    mkSubst (dom dec.bounds) gs = some sigma ->
    Step d (.field (.new c gs (vs1 ++ vs)) f)
      (.field (.new dec.sup (Ty.substList sigma dec.supArgs) vs) f)
\end{lstlisting}
\end{leanside}
\mapping{$[\Delta=\pl{G}]$ is \texttt{sigma}.}
\end{cmp}

\begin{cmp}{\textsc{innerFieldSuper}}
\begin{paperside}
\fit{$\inferrule{ }{
    \step{\innerobj{C}{\pl{G}}{C'}{\pl{G}'}{\pl{v}}{\pl{M}}\fstop f}{\obj{C'}{\pl{G}'}{\pl{v}}\fstop f}
}$}
\end{paperside}
\begin{leanside}
\begin{lstlisting}[style=leanside]
| innerFieldSuper {c fs c' gs vs ms f} :
    (Expr.newLocal c fs c' gs vs ms).isValue = true ->
    Step d (.field (.newLocal c fs c' gs vs ms) f) (.field (.new c' gs vs) f)
\end{lstlisting}
\end{leanside}
\mapping{The premise \texttt{isValue} is the grammar of $v$ and $G$.}
\end{cmp}

\begin{cmp}{\textsc{ifThen}}
\begin{paperside}
\fit{$\inferrule{
    \emptyset\vdash \mathsf{typeOf}(v) <: \gen{C'}{\_}\\
    \mathsf{fresh}(e_1,e'_1)
}{
    \step{\ifthen{v}{\gen{C'}{\pl{T}}}{x}{e_1}{e_2}}{\tsubst{e'_1}{x}{v}}
}$}
\end{paperside}
\begin{leanside}
\begin{lstlisting}[style=leanside]
| ifThen {v c' ts ts' tv x e1 e1' e2} :
    v.isValue = true -> typeOf v = some tv ->
    Sub d [] tv (.cls c' ts') ->
    Fresh (avoidNames d (.ifInst v (.cls c' ts) x e1 e2)) e1 e1' ->
    Step d (.ifInst v (.cls c' ts) x e1 e2) (e1'.vsubst [(x, v)])
\end{lstlisting}
\end{leanside}
\mapping{The $\_$ in $\gen{C'}{\_}$ is the implicit \texttt{ts'}.}
\end{cmp}

\begin{cmp}{\textsc{ifElse}}
\begin{paperside}
\fit{$\inferrule{
    \emptyset\not\vdash \mathsf{typeOf}(v) <: \gen{C'}{\_}
}{
    \step{\ifthen{v}{\gen{C'}{\pl{T}}}{x}{e_1}{e_2}}{e_2}
}$}
\end{paperside}
\begin{leanside}
\begin{lstlisting}[style=leanside]
| ifElse {v c' ts tv x e1 e2} :
    v.isValue = true -> typeOf v = some tv ->
    (forall ts', Not (Sub d [] tv (.cls c' ts'))) ->
    Step d (.ifInst v (.cls c' ts) x e1 e2) e2
\end{lstlisting}
\end{leanside}
\mapping{$\not\vdash\ldots<:\gen{C'}{\_}$: no type arguments make the subtyping hold.}
\end{cmp}

\begin{cmp}{\textsc{simpleAddSuper}}
\begin{paperside}
\fit{$\inferrule{
    v = \obj{C}{\pl{G}}{\_}
}{
    \step{\call{v}{\pl{G}'}{m}{\pl{v}}}{\explicitcall{v}{\gen{C}{\pl{G}}}{\pl{G}'}{m}{\pl{v}}}
}$}
\end{paperside}
\begin{leanside}
\begin{lstlisting}[style=leanside]
| simpleAddSuper {c gs vs0 gs' m vs} :
    (Expr.new c gs vs0).isValue = true -> Ty.allGround gs' = true ->
    Expr.allValues vs = true ->
    Step d (.call (.new c gs vs0) gs' m vs)
      (.superCall (.new c gs vs0) (.cls c gs) gs' m vs)
\end{lstlisting}
\end{leanside}
\mapping{The premise $v=\obj{C}{\pl{G}}{\_}$ is the shape of the receiver in the conclusion.}
\end{cmp}

\begin{cmp}{\textsc{simpleFound}}
\begin{paperside}
\fit{$\inferrule{
    \declared{\subclass{C}{\Delta}{\_}{\_}{\pl{M}}{\_}}\\
    M\in\pl{M}\\\\
    \mathsf{fresh}(M,\method{\Delta'}{T_0}{m}{T_1\gap x_1,\ldots,T_n\gap x_n}{\ret{e_0}})\\\\
    e = \tsubst{\tsubst{e_0}{\Delta,\Delta'}{\pl{G},\pl{G}'}}{\this,x_1,\ldots,x_n}{v_0,\ldots,v_n}
}{
    \step{\explicitcall{v_0}{\gen{C}{\pl{G}}}{\pl{G}'}{m}{v_1,\ldots,v_n}}{e}
}$}
\end{paperside}
\begin{leanside}
\begin{lstlisting}[style=leanside]
| simpleFound {v0 c gs gs' m vs mt delta' t0 ps e0 sigma} {dec : Dec} :
    List.Mem dec d -> dec.name = c ->
    v0.isValue = true -> Ty.allGround gs = true -> Ty.allGround gs' = true ->
    Expr.allValues vs = true ->
    List.Mem mt dec.methods ->
    FreshM (avoidNames d (.superCall v0 (.cls c gs) gs' m vs)) mt
      (.mk delta' t0 m ps e0) ->
    mkSubst (dom dec.bounds ++ dom delta') (gs ++ gs') = some sigma ->
    ps.length = vs.length ->
    Step d (.superCall v0 (.cls c gs) gs' m vs)
      ((e0.tsubst sigma).vsubst ((thisName, v0) :: (ps.map Prod.snd).zip vs))
\end{lstlisting}
\end{leanside}
\mapping{$[\Delta,\Delta'=\pl{G},\pl{G}']$ is \texttt{sigma}; $x_1,\ldots,x_n$ is \texttt{ps.map Prod.snd}, as many as the arguments \texttt{vs}.}
\end{cmp}

\begin{cmp}{\textsc{innerFound}}
\begin{paperside}
\fit{$\inferrule{
    v_0 = \innerobj{\_}{\_}{\_}{\_}{\_}{\pl{M}}\\
    M\in\pl{M}\\\\
    \mathsf{fresh}(M,\method{\Delta'}{T_0}{m}{T_1\gap x_1,\ldots,T_n\gap x_n}{\ret{e_0}})\\\\
    e = \tsubst{\tsubst{e_0}{\Delta'}{\pl{G}'}}{\this,x_1,\ldots,x_n}{v_0,\ldots,v_n}
}{
    \step{\call{v_0}{\pl{G}'}{m}{v_1,\ldots,v_n}}{e}
}$}
\end{paperside}
\begin{leanside}
\begin{lstlisting}[style=leanside]
| innerFound {c fs c' gs vs0 ms gs' m vs mt delta' t0 ps e0 sigma} :
    (Expr.newLocal c fs c' gs vs0 ms).isValue = true -> Ty.allGround gs' = true ->
    Expr.allValues vs = true ->
    List.Mem mt ms ->
    FreshM (avoidNames d (.call (.newLocal c fs c' gs vs0 ms) gs' m vs)) mt
      (.mk delta' t0 m ps e0) ->
    mkSubst (dom delta') gs' = some sigma ->
    ps.length = vs.length ->
    Step d (.call (.newLocal c fs c' gs vs0 ms) gs' m vs)
      ((e0.tsubst sigma).vsubst
        ((thisName, .newLocal c fs c' gs vs0 ms) :: (ps.map Prod.snd).zip vs))
\end{lstlisting}
\end{leanside}
\mapping{As \textsc{simpleFound}, with $v_0$ written out as \texttt{.newLocal c fs c' gs vs0 ms}.}
\end{cmp}

\begin{cmp}{\textsc{innerSuper}}
\begin{paperside}
\fit{$\inferrule{
    v_0 = \innerobj{\_}{\_}{C}{\pl{G}}{\_}{\pl{M}}\\
    m\not\in \pl{M}
}{
    \step{\call{v_0}{\pl{G}'}{m}{v_1,\ldots,v_n}}{\explicitcall{v_0}{\gen{C}{\pl{G}}}{\pl{G}'}{m}{v_1,\ldots,v_n}}
}$}
\end{paperside}
\begin{leanside}
\begin{lstlisting}[style=leanside]
| innerSuper {c fs c' gs vs0 ms gs' m vs} :
    (Expr.newLocal c fs c' gs vs0 ms).isValue = true -> Ty.allGround gs' = true ->
    Expr.allValues vs = true ->
    (ms.map Method.name).contains m = false ->
    Step d (.call (.newLocal c fs c' gs vs0 ms) gs' m vs)
      (.superCall (.newLocal c fs c' gs vs0 ms) (.cls c' gs) gs' m vs)
\end{lstlisting}
\end{leanside}
\mapping{The superclass $C$ of the paper is \texttt{c'} in Lean.}
\end{cmp}

\begin{cmp}{\textsc{searchSuper}}
\begin{paperside}
\fit{$\inferrule{
    \declared{\subclass{C}{\Delta}{C'}{\pl{T}}{\pl{M}}{\_}}\\
    m\not\in \pl{M}\\
    C\neq\Object
}{
    \step{\explicitcall{v}{\gen{C}{\pl{G}}}{\pl{G}'}{m}{\pl{v}}}{\explicitcall{v}{\gen{C'}{\tsubst{\pl{T}}{\Delta}{\pl{G}}}}{\pl{G}'}{m}{\pl{v}}}
}$}
\end{paperside}
\begin{leanside}
\begin{lstlisting}[style=leanside]
| searchSuper {v c gs gs' m vs sigma} {dec : Dec} :
    List.Mem dec d -> dec.name = c ->
    v.isValue = true -> Ty.allGround gs = true -> Ty.allGround gs' = true ->
    Expr.allValues vs = true ->
    (dec.methods.map Method.name).contains m = false ->
    Not (c = objectName) ->
    mkSubst (dom dec.bounds) gs = some sigma ->
    Step d (.superCall v (.cls c gs) gs' m vs)
      (.superCall v (.cls dec.sup (Ty.substList sigma dec.supArgs)) gs' m vs)
\end{lstlisting}
\end{leanside}
\mapping{$m\not\in\pl{M}$ is \texttt{contains m = false} on the method names.}
\end{cmp}

\begin{cmp}{\textsc{ctx} and $\rightarrow^*$}
\begin{paperside}
\fit{$\inferrule{
    \step{e}{e'}
}{
    \step{\ctx{e}}{\ctx{e'}}
}$}\\[1ex]
$e_0\rightarrow^* e$: zero or more steps.
\end{paperside}
\begin{leanside}
\begin{lstlisting}[style=leanside]
| ctx {E : ECtx} {e e'} :
    E.ok = true -> Step d e e' -> Step d (E.plug e) (E.plug e')
inductive Steps (d : Decls) : Expr -> Expr -> Prop
  | refl {e} : Steps d e e
  | step {e e' e''} : Step d e e' -> Steps d e' e'' -> Steps d e e''
\end{lstlisting}
\end{leanside}
\mapping{$\mathcal{E}$ is \texttt{E} with \texttt{E.ok}; $\rightarrow^*$ is \texttt{Steps d}.}
\end{cmp}

\subsubsection{Subtyping and Kinding}

\begin{cmp}{\textsc{sub-refl}, \textsc{sub-var}}
\begin{paperside}
\fit{$\inferrule{ }{\subtype{\Delta}{T}{T}}$}\qquad
\fit{$\inferrule{
  \ext{X}{T}\in\Delta\\
  \subtype{\Delta}{T}{T'}
}{
  \subtype{\Delta}{X}{T'}
}$}
\end{paperside}
\begin{leanside}
\begin{lstlisting}[style=leanside]
inductive Sub (d : Decls) (delta : Bounds) : Ty -> Ty -> Prop
  | refl {t} : Sub d delta t t
  | var {x t t'} :
      delta.lookup x = some t -> Sub d delta t t' -> Sub d delta (.var x) t'
\end{lstlisting}
\end{leanside}
\mapping{$\ext{X}{T}\in\Delta$ is \texttt{delta.lookup x = some t}; by A5 there is one bound per $X$.}
\end{cmp}

\begin{cmp}{\textsc{sub-inherit}}
\begin{paperside}
\fit{$\inferrule{
  \declared{\subclass{C}{\Delta_0}{C'}{\pl{T}}{\_}{\_}}\\\\
  \subtype{\Delta}{\tsubst{\gen{C'}{\pl{T}}}{\Delta_0}{\pl{T}_0}}{T'}
}{
  \subtype{\Delta}{\gen{C}{\pl{T}_0}}{T'}
}$}
\end{paperside}
\begin{leanside}
\begin{lstlisting}[style=leanside]
  | inherit {c ts0 t' sigma} {dec : Dec} :
      List.Mem dec d -> dec.name = c ->
      mkSubst (dom dec.bounds) ts0 = some sigma ->
      Sub d delta ((Ty.cls dec.sup dec.supArgs).subst sigma) t' ->
      Sub d delta (.cls c ts0) t'
\end{lstlisting}
\end{leanside}
\mapping{$\Delta_0$ is \texttt{dec.bounds}; $\gen{C'}{\pl{T}}$ is \texttt{Ty.cls dec.sup dec.supArgs}.}
\end{cmp}

\begin{cmp}{\textsc{sub-inherit-local}}
\begin{paperside}
\fit{$\inferrule{
  \declared{\innerobj{C}{\pl{F}}{C'}{\pl{T}'}{\_}{\_}}\\\\
  \subtype{\Delta}{\tsubst{\gen{C'}{\pl{T}'}}{\pl{F}}{\pl{T}}}{T'}
}{
  \subtype{\Delta}{\gen{C}{\pl{T}}}{T'}
}$}
\end{paperside}
\begin{leanside}
\begin{lstlisting}[style=leanside]
  | inheritLocal {c fs c' ts' es ms ts t' sigma} :
      Declared d (.loc c fs c' ts' es ms) ->
      mkSubstF fs ts = some sigma ->
      Sub d delta ((Ty.cls c' ts').subst sigma) t' ->
      Sub d delta (.cls c ts) t'
\end{lstlisting}
\end{leanside}
\mapping{$[\pl{F}=\pl{T}]$ is \texttt{mkSubstF fs ts}.}
\end{cmp}

\begin{cmp}{\textsc{wf-var}, \textsc{wf-class}}
\begin{paperside}
\fit{$\inferrule{
  \ext{X}{\_}\in\Delta
}{
  \kind{\Delta}{X}
}$}\\[1ex]
\fit{$\inferrule{
  \mathsf{scope}(\gen{C}{\pl{T}})=\Delta_0;\pl{T}_0\\
  \kind{\Delta}{\pl{T}}\\\\
  \subtype{\Delta}{\pl{T}_0}{\tsubst{\mathrm{cod}(\Delta_0)}{\Delta_0}{\pl{T}_0}}
}{
  \kind{\Delta}{\gen{C}{\pl{T}}}
}$}
\end{paperside}
\begin{leanside}
\begin{lstlisting}[style=leanside]
inductive Kind (d : Decls) (delta : Bounds) : Ty -> Prop
  | var {x} : (delta.lookup x).isSome = true -> Kind d delta (.var x)
  | cls {c ts delta0 ts0 sigma} :
      scope d c ts = some (delta0, ts0) ->
      (forall t, List.Mem t ts -> Kind d delta t) ->
      mkSubst (dom delta0) ts0 = some sigma ->
      (forall p, List.Mem p (ts0.zip (Ty.substList sigma (cod delta0))) ->
        Sub d delta p.1 p.2) ->
      Kind d delta (.cls c ts)
\end{lstlisting}
\end{leanside}
\mapping{The subtyping on sequences is one \texttt{Sub} for each pair of \texttt{zip}.}
\end{cmp}

\subsubsection{Checked Cast}

\begin{cmp}{\textsc{cast-up}, \textsc{cast-var}}
\begin{paperside}
\fit{$\inferrule{
  \subtype{\Delta}{T}{T'}
}{
  \cast{\Delta}{T}{T'}
}$}\qquad
\fit{$\inferrule{
  \ext{X}{T}\in\Delta\\\\
  \cast{\Delta}{T}{T'}
}{
  \cast{\Delta}{X}{T'}
}$}
\end{paperside}
\begin{leanside}
\begin{lstlisting}[style=leanside]
inductive Cast (d : Decls) (delta : Bounds) : Ty -> Ty -> Prop
  | up {t t'} : Sub d delta t t' -> Cast d delta t t'
  | var {x t t'} :
      delta.lookup x = some t -> Cast d delta t t' -> Cast d delta (.var x) t'
\end{lstlisting}
\end{leanside}
\mapping{$\cast{\Delta}{T}{T'}$ is \texttt{Cast d delta T T'}.}
\end{cmp}

\begin{cmp}{\textsc{cast-down}}
\begin{paperside}
\fit{$\inferrule{
  \subtype{\Delta}{\gen{C}{\pl{T}}}{\gen{C'}{\pl{T}_0}}\\
  \generic{C}=\gen{C}{\pl{X}}\\\\
  \subtype{\emptyset}{\gen{C}{\pl{X}}}{\gen{C'}{\pl{T}'}}\\
  \pl{X}\subseteq\ftv{\pl{T}'}
}{
  \cast{\Delta}{\gen{C'}{\pl{T}_0}}{\gen{C}{\pl{T}}}
}$}
\end{paperside}
\begin{leanside}
\begin{lstlisting}[style=leanside]
  | down {c ts c' ts0 xs ts'} :
      Sub d delta (.cls c ts) (.cls c' ts0) ->
      generic d c = some xs ->
      Sub d [] (.cls c (vars xs)) (.cls c' ts') ->
      (forall x, List.Mem x xs -> List.Mem x (Ty.ftvList ts')) ->
      Cast d delta (.cls c' ts0) (.cls c ts)
\end{lstlisting}
\end{leanside}
\mapping{$\pl{X}$ is \texttt{xs}; $\gen{C}{\pl{X}}$ is \texttt{.cls c (vars xs)}.}
\end{cmp}

\subsubsection{Typing}

\begin{cmp}{\textsc{subs}, \textsc{var}, \textsc{new-object}}
\begin{paperside}
\fit{$\inferrule{
	\typing{\Delta}{\Gamma}{e}{T_1}\\
	\kind{\Delta}{T_2}\\
	\subtype{\Delta}{T_1}{T_2}
}{
	\typing{\Delta}{\Gamma}{e}{T_2}
}$}\\[1ex]
\fit{$\inferrule{ }{\typing{\Delta}{\Gamma}{x}{\Gamma(x)}}$}\qquad
\fit{$\inferrule{ }{
  \typing{\Delta}{\Gamma}{\obj{\Object}{}{}}{\gen{\Object}{}}
}$}
\end{paperside}
\begin{leanside}
\begin{lstlisting}[style=leanside]
inductive Typing (d : Decls) (java : Bool) : Bounds -> Ctx -> Expr -> Ty -> Prop
  | subs {delta gamma e t1 t2} :
      Typing d java delta gamma e t1 -> Kind d delta t2 -> Sub d delta t1 t2 ->
      Typing d java delta gamma e t2
  | var {delta gamma x t} :
      gamma.lookup x = some t -> Typing d java delta gamma (.var x) t
  | newObject {delta gamma} :
      Typing d java delta gamma (.new objectName [] []) (.cls objectName [])
\end{lstlisting}
\end{leanside}
\mapping{The flag \texttt{java} selects the variant: \texttt{false} is $\zeta(\mathsf{funnelled})$, \texttt{true} is $\zeta(\mathsf{Java})$. Only \textsc{new-local} and \textsc{new-local-java} read it.}
\end{cmp}

\begin{cmp}{\textsc{declField}}
\begin{paperside}
\fit{$\inferrule{
  \typing{\Delta}{\Gamma}{e}{\classtype{C}{\pl{T}'}}\\\\
  \declared{\subclass{C}{\Delta'}{\_}{\_}{\_}{\pl{S}}}\\\\
  T\gap f\semi\in\pl{S}
}{
  \typing{\Delta}{\Gamma}{e\fstop f}{\tsubst{T}{\Delta'}{\pl{T}'}}
}$}
\end{paperside}
\begin{leanside}
\begin{lstlisting}[style=leanside]
  | declField {delta gamma e c ts' t f sigma} {dec : Dec} :
      Typing d java delta gamma e (.cls c ts') ->
      List.Mem dec d -> dec.name = c ->
      List.Mem (t, f) dec.fields ->
      mkSubst (dom dec.bounds) ts' = some sigma ->
      Typing d java delta gamma (.field e f) (t.subst sigma)
\end{lstlisting}
\end{leanside}
\mapping{$\Delta'$ is \texttt{dec.bounds}; $\pl{S}$ is \texttt{dec.fields}.}
\end{cmp}

\begin{cmp}{\textsc{call}}
\begin{paperside}
\fit{$\inferrule{
  \typing{\Delta}{\Gamma}{e_0}{\classtype{C_0}{\pl{T}_0}}\\\\
  \mathsf{scope}(\gen{C_0}{\pl{T}_0})=\Delta_0;\pl{T}_2\\\\
  \method{\Delta_1}{T_0}{m}{T_1\gap x_1,\ldots,T_n\gap x_n}{\_} \in C_0\\\\
  \forall i\in 0..n\ldotp T'_i=\tsubst{T_i}{\Delta_0,\Delta_1}{\pl{T}_2,\pl{T}_1}\\\\
  \kind{\Delta}{\pl{T}_1}\\\\
  \subtype{\Delta}{\pl{T}_1}{\tsubst{\mathrm{cod}(\Delta_1)}{\Delta_0,\Delta_1}{\pl{T}_2,\pl{T}_1}}\\\\
  \typing{\Delta}{\Gamma}{e_1,\ldots,e_n}{T'_1,\ldots,T'_n}
}{
  \typing{\Delta}{\Gamma}{\call{e_0}{\pl{T}_1}{m}{e_1,\ldots,e_n}}{T'_0}
}$}
\end{paperside}
\begin{leanside}
\begin{lstlisting}[style=leanside]
  | call {delta gamma e0 c0 ts0 delta0 ts2 delta1 t0 m ps body ts1 es sigma} :
      Typing d java delta gamma e0 (.cls c0 ts0) ->
      scope d c0 ts0 = some (delta0, ts2) ->
      MethodIn d (.mk delta1 t0 m ps body) c0 ->
      mkSubst (dom delta0 ++ dom delta1) (ts2 ++ ts1) = some sigma ->
      (forall t, List.Mem t ts1 -> Kind d delta t) ->
      (forall p, List.Mem p (ts1.zip (Ty.substList sigma (cod delta1))) ->
        Sub d delta p.1 p.2) ->
      es.length = ps.length ->
      (forall p, List.Mem p (es.zip (Ty.substList sigma (ps.map Prod.fst))) ->
        Typing d java delta gamma p.1 p.2) ->
      Typing d java delta gamma (.call e0 ts1 m es) (t0.subst sigma)
\end{lstlisting}
\end{leanside}
\mapping{$[\Delta_0,\Delta_1=\pl{T}_2,\pl{T}_1]$ is \texttt{sigma}; $T'_i$ is written out as \texttt{subst sigma}.}
\end{cmp}

\begin{cmp}{\textsc{call-lb-runtime}}
\begin{paperside}
\fit{$\inferrule{
  \subtype{\Delta}{\gen{C}{\pl{T}}}{\gen{C_0}{\pl{T}_0}}\\\\
  \declared{\subclass{C_0}{\Delta_0}{\_}{\_}{\_}{\_}}\\\\
  \method{\Delta_1}{T_0}{m}{T_1\gap x_1,\ldots,T_n\gap x_n}{\_} \in C_0\\\\
  \forall i\in 0..n\ldotp T'_i=\tsubst{T_i}{\Delta_0,\Delta_1}{\pl{T}_0,\pl{T}_1}\\\\
  \declared{\subclass{C}{\_}{\_}{\_}{\_}{\_}}\\\\
  \kind{\Delta}{\gen{C}{\pl{T}}}\\
  \kind{\Delta}{\pl{T}_1}\\\\
  \subtype{\Delta}{\pl{T}_1}{\tsubst{\mathrm{cod}(\Delta_1)}{\Delta_0,\Delta_1}{\pl{T}_0,\pl{T}_1}}\\\\
  \typing{\Delta}{\Gamma}{e_0,\ldots,e_n}{\gen{C}{\pl{T}},T'_1,\ldots,T'_n}
}{
  \typing{\Delta}{\Gamma}{
    \explicitcall{e_0}{\gen{C}{\pl{T}}}{\pl{T}_1}{m}{e_1,\ldots,e_n}
  }{T'_0}
}$}
\end{paperside}
\begin{leanside}
\begin{lstlisting}[style=leanside]
  | superCall {delta gamma e0 c ts c0 ts0 delta1 t0 m ps body ts1 es sigma}
      {dec0 dec : Dec} :
      Sub d delta (.cls c ts) (.cls c0 ts0) ->
      List.Mem dec0 d -> dec0.name = c0 ->
      MethodIn d (.mk delta1 t0 m ps body) c0 ->
      mkSubst (dom dec0.bounds ++ dom delta1) (ts0 ++ ts1) = some sigma ->
      List.Mem dec d -> dec.name = c ->
      Kind d delta (.cls c ts) ->
      (forall t, List.Mem t ts1 -> Kind d delta t) ->
      (forall p, List.Mem p (ts1.zip (Ty.substList sigma (cod delta1))) ->
        Sub d delta p.1 p.2) ->
      Typing d java delta gamma e0 (.cls c ts) ->
      es.length = ps.length ->
      (forall p, List.Mem p (es.zip (Ty.substList sigma (ps.map Prod.fst))) ->
        Typing d java delta gamma p.1 p.2) ->
      Typing d java delta gamma (.superCall e0 (.cls c ts) ts1 m es) (t0.subst sigma)
\end{lstlisting}
\end{leanside}
\mapping{$\Delta_0$ is \texttt{dec0.bounds}; the premise on $e_0$ is the first element of the paper's sequence.}
\end{cmp}

\begin{cmp}{\textsc{new-top}}
\begin{paperside}
\fit{$\inferrule{
  \declared{\subclass{C}{\Delta'}{C'}{\pl{T}}{\_}{T_1\gap f_1\semi\ldots T_n\gap f_n\semi}}\\
  \kind{\Delta}{\gen{C}{\pl{T}'}}\\\\
  \typing{\Delta}{\Gamma}{e_1,\ldots,e_n}{\tsubst{T_1,\ldots,T_n}{\Delta'}{\pl{T}'}}\\
  \typing{\Delta}{\Gamma}{\obj{C'}{\tsubst{\pl{T}}{\Delta'}{\pl{T}'}}{\pl{e}}}{\_}
}{
  \typing{\Delta}{\Gamma}{\obj{C}{\pl{T}'}{e_1,\ldots,e_n,\pl{e}}}{\gen{C}{\pl{T}'}}
}$}
\end{paperside}
\begin{leanside}
\begin{lstlisting}[style=leanside]
  | newTop {delta gamma c ts' es1 es2 sigma t''} {dec : Dec} :
      List.Mem dec d -> dec.name = c ->
      Kind d delta (.cls c ts') ->
      mkSubst (dom dec.bounds) ts' = some sigma ->
      es1.length = dec.fields.length ->
      (forall p, List.Mem p (es1.zip (Ty.substList sigma (dec.fields.map Prod.fst))) ->
        Typing d java delta gamma p.1 p.2) ->
      Typing d java delta gamma (.new dec.sup (Ty.substList sigma dec.supArgs) es2) t'' ->
      Typing d java delta gamma (.new c ts' (es1 ++ es2)) (.cls c ts')
\end{lstlisting}
\end{leanside}
\mapping{$e_1,\ldots,e_n$ is \texttt{es1}, as long as the fields; $\pl{e}$ is \texttt{es2}; the $\_$ type is \texttt{t''}.}
\end{cmp}

\begin{cmp}{\textsc{new-local}}
\begin{paperside}
\fit{$\inferrule{
  \innerobj{C}{\pl{F}}{C_0}{\pl{T}}{\pl{e}}{\pl{M}}\gap\mathsf{consistent}\\\\
  \pairing{\Delta}{\pl{F}}{|\pl{F}|}=\Delta';\_\\
  \ftv{\mathrm{cod}(\Delta')}\subseteq\mathrm{dom}(\Delta')\\\\
  \kind{\Delta'}{\gen{C}{|\pl{F}|}}\\
  \kind{\Delta'}{\gen{C_0}{\pl{T}}}\\
  \typing{\Delta}{\Gamma}{\obj{C_0}{\pl{T}}{\pl{e}}}{\_}\\\\
  \mok{\Delta'}{\Gamma|_{\ftv{|\pl{F}|}}\setminus\this,\gen{C}{|\pl{F}|}\gap\this}{\pl{M}}\\\\
  \mathsf{overrideOk}(\Delta',\gen{C_0}{\pl{T}},\pl{M})
}{
  \typing{\Delta}{\Gamma}{
    \innerobj{C}{\pl{F}}{C_0}{\pl{T}}{\pl{e}}{\pl{M}}
  }{\gen{C}{|\pl{F}|}}
}$}\\[1ex]
For a set of type variables $\pl{X}$, $\Gamma|_{\pl{X}}$ keeps the variables $T\gap x$ of $\Gamma$ such that $\ftv{T}\subseteq\pl{X}$.
$\Gamma\setminus x$ is $\Gamma$ without the variable $x$.
\end{paperside}
\begin{leanside}
\begin{lstlisting}[style=leanside]
  | newLocal {delta gamma c fs c0 ts es ms delta' ts'' t''} :
      java = false ->
      Consistent d (.loc c fs c0 ts es ms) ->
      pairing delta fs (fs.map Funnel.abs) = some (delta', ts'') ->
      (forall x, List.Mem x (Ty.ftvList (cod delta')) -> List.Mem x (dom delta')) ->
      Kind d delta' (.cls c (fs.map Funnel.abs)) ->
      Kind d delta' (.cls c0 ts) ->
      Typing d java delta gamma (.new c0 ts es) t'' ->
      (forall mt, List.Mem mt ms ->
        MethodOk d java delta'
          (restrict gamma (Ty.ftvList (fs.map Funnel.abs)) ++
            [(thisName, .cls c (fs.map Funnel.abs))]) mt) ->
      OverrideOk d delta' (.cls c0 ts) ms ->
      Typing d java delta gamma (.newLocal c fs c0 ts es ms) (.cls c (fs.map Funnel.abs))
def restrict (gamma : Ctx) (xs : List Name) : Ctx :=
  gamma.filter fun p => p.1 != thisName && p.2.ftv.all fun x => xs.contains x
\end{lstlisting}
\end{leanside}
\mapping{$|\pl{F}|$ is \texttt{fs.map Funnel.abs}; $\Gamma|_{\ftv{|\pl{F}|}}\setminus\this$ is \texttt{restrict}; adding $\gen{C}{|\pl{F}|}\gap\this$ appends a pair, and \texttt{lookup} finds it since \texttt{restrict} removed the old \Q@this@.}
\end{cmp}

\begin{cmp}{\textsc{new-local-java}}
\begin{paperside}
\fit{$\inferrule{
  \kind{\Delta}{\gen{C'}{\pl{T}}}\\
  \typing{\Delta}{\Gamma}{\obj{C'}{\pl{T}}{\pl{e}}}{\_}\\\\
  \mok{\Delta}{\Gamma\setminus\this,\gen{C}{}\gap\this}{\pl{M}}\\\\
  \mathsf{overrideOk}(\Delta,\gen{C'}{\pl{T}},\pl{M})
}{
  \typing{\Delta}{\Gamma}{
    \innerobj{C}{}{C'}{\pl{T}}{\pl{e}}{\pl{M}}
  }{\gen{C}{}}
}$}
\end{paperside}
\begin{leanside}
\begin{lstlisting}[style=leanside]
  | newLocalJava {delta gamma c c0 ts es ms t''} :
      java = true ->
      Kind d delta (.cls c0 ts) ->
      Typing d java delta gamma (.new c0 ts es) t'' ->
      (forall mt, List.Mem mt ms ->
        MethodOk d java delta (withoutThis gamma ++ [(thisName, .cls c [])]) mt) ->
      OverrideOk d delta (.cls c0 ts) ms ->
      Typing d java delta gamma (.newLocal c [] c0 ts es ms) (.cls c [])
def withoutThis (gamma : Ctx) : Ctx := gamma.filter fun p => p.1 != thisName
\end{lstlisting}
\end{leanside}
\mapping{$C'$ of the paper is \texttt{c0} in Lean.}
\end{cmp}

\begin{cmp}{\textsc{if-instanceof}}
\begin{paperside}
\fit{$\inferrule{
  \typing{\Delta}{\Gamma}{e_0}{T_0}\\
  \kind{\Delta}{T}\\
  \cast{\Delta}{T_0}{T}\\\\
  \typing{\Delta}{\Gamma,T\gap x}{e_1}{T'}\\
  \typing{\Delta}{\Gamma}{e_2}{T'}
}{
  \typing{\Delta}{\Gamma}{\ifthen{e_0}{T}{x}{e_1}{e_2}}{T'}
}$}
\end{paperside}
\begin{leanside}
\begin{lstlisting}[style=leanside]
  | ifInst {delta gamma e0 t0 t x e1 e2 t'} :
      Typing d java delta gamma e0 t0 -> Kind d delta t -> Cast d delta t0 t ->
      Typing d java delta (gamma ++ [(x, t)]) e1 t' ->
      Typing d java delta gamma e2 t' ->
      Typing d java delta gamma (.ifInst e0 t x e1 e2) t'
\end{lstlisting}
\end{leanside}
\mapping{$\Gamma,T\gap x$ is \texttt{gamma ++ [(x, t)]}; by A7, $x$ is not already in $\Gamma$.}
\end{cmp}

\subsubsection{Consistency and Overriding}

\begin{cmp}{$K\gap\mathsf{consistent}$}
\begin{paperside}
\fit{$\begin{array}{@{}l@{}}
\innerobj{C}{\pl{F}}{C_0}{\pl{T}}{\pl{e}}{\pl{M}}\gap\mathsf{consistent}\\
\quad\textrm{where }\declared{\innerobj{C}{\pl{F}'}{C_0}{\pl{T}'}{\_}{\pl{M}'}}\\
\quad\phantom{\textrm{where }}\pl{T}=\tsubst{\pl{T}'}{\pl{F}'}{|\pl{F}|}\\
\quad\phantom{\textrm{where }}\mathrm{dom}(\pl{M}')=\mathrm{dom}(\pl{M})\\
\quad\phantom{\textrm{where }}\mathsf{overrideOk}(\gen{C}{|\pl{F}|},\pl{M})
\end{array}$}\\[1ex]
$\mathrm{dom}(\pl{M})$ is the set of the names of the methods in $\pl{M}$.
\end{paperside}
\begin{leanside}
\begin{lstlisting}[style=leanside]
def Consistent (d : Decls) : K -> Prop
  | .loc c fs c0 ts _ ms =>
    exists fs' ts' es' ms' sigma,
      Declared d (.loc c fs' c0 ts' es' ms') /\
      mkSubstF fs' (fs.map Funnel.abs) = some sigma /\
      ts = Ty.substList sigma ts' /\
      (forall m, List.Mem m (ms'.map Method.name) <->
        List.Mem m (ms.map Method.name)) /\
      forall mt, List.Mem mt ms -> OverrideOk1 d (.cls c (fs.map Funnel.abs)) mt
  | .top _ => False
\end{lstlisting}
\end{leanside}
\mapping{$\mathsf{overrideOk}(\gen{C}{|\pl{F}|},\pl{M})$ on the sequence $\pl{M}$ is one \texttt{OverrideOk1} per method.}
\end{cmp}

\begin{cmp}{$\mathsf{overrideOk}(\Delta,T,\pl{M})$}
\begin{paperside}
\fit{$\begin{array}{@{}l@{}}
\mathsf{overrideOk}(\Delta,T,\pl{M})\textrm{ iff }\mathsf{overrideOk}(\gen{C}{\pl{T}},M)\\
\quad\textrm{for every }M\in\pl{M}\textrm{ and every }\gen{C}{\pl{T}}\textrm{ such that }\subtype{\Delta}{T}{\gen{C}{\pl{T}}}
\end{array}$}
\end{paperside}
\begin{leanside}
\begin{lstlisting}[style=leanside]
def OverrideOk (d : Decls) (delta : Bounds) (t : Ty) (ms : List Method) : Prop :=
  forall mt c ts, List.Mem mt ms -> Sub d delta t (.cls c ts) ->
    OverrideOk1 d (.cls c ts) mt
\end{lstlisting}
\end{leanside}
\mapping{$\mathsf{overrideOk}(T,M)$ is \texttt{OverrideOk1}.}
\end{cmp}

\begin{cmp}{$\mathsf{overrideOk}(T,M)$}
\begin{paperside}
\fit{$\begin{array}{@{}l@{}}
\mathsf{overrideOk}(\gen{C}{\pl{T}},\method{\_}{\_}{m}{\_}{\_})\textrm{ iff }m\not\in\pl{M}\\
\quad\textrm{where }\declared{\subclass{C}{\_}{\_}{\_}{\pl{M}}{\_}}\\[1ex]
\mathsf{overrideOk}(\gen{C}{\pl{T}},\method{\Delta'}{T'_0}{m}{T'_1\gap x_1,\ldots,T'_n\gap x_n}{\_})\\
\quad\textrm{where }\declared{\subclass{C}{\Delta}{\_}{\_}{\pl{M}}{\_}}\\
\quad\phantom{\textrm{where }}\method{\Delta''}{T_0}{m}{T_1\gap x'_1,\ldots,T_n\gap x'_n}{\_}\in\pl{M}\\
\quad\phantom{\textrm{where }}\tsubst{\mathrm{cod}(\Delta'')}{\Delta,\Delta''}{\pl{T},\mathrm{dom}(\Delta')}=\mathrm{cod}(\Delta')\\
\quad\phantom{\textrm{where }}\tsubst{T_0,\ldots,T_n}{\Delta,\Delta''}{\pl{T},\mathrm{dom}(\Delta')}=T'_0,\ldots,T'_n\\[1ex]
\mathsf{overrideOk}(\gen{C}{\pl{T}},\method{\Delta'}{T'_0}{m}{T'_1\gap x_1,\ldots,T'_n\gap x_n}{\_})\\
\quad\textrm{where }\declared{\innerobj{C}{\pl{F}}{\_}{\_}{\_}{\pl{M}}}\\
\quad\phantom{\textrm{where }}\method{\Delta''}{T_0}{m}{T_1\gap x'_1,\ldots,T_n\gap x'_n}{\_}\in\pl{M}\\
\quad\phantom{\textrm{where }}\tsubst{\mathrm{cod}(\Delta'')}{\pl{F},\Delta''}{\pl{T},\mathrm{dom}(\Delta')}=\mathrm{cod}(\Delta')\\
\quad\phantom{\textrm{where }}\tsubst{T_0,\ldots,T_n}{\pl{F},\Delta''}{\pl{T},\mathrm{dom}(\Delta')}=T'_0,\ldots,T'_n
\end{array}$}
\end{paperside}
\begin{leanside}
\begin{lstlisting}[style=leanside]
def OverrideOk1 (d : Decls) (t : Ty) (mt : Method) : Prop :=
  match t, mt with
  | .cls c ts, .mk delta' t0' m ps' _ =>
    (exists dec : Dec, List.Mem dec d /\ dec.name = c /\
       Not (List.Mem m (dec.methods.map Method.name)))
    \/ (exists (dec : Dec) (delta'' : Bounds) (t0 : Ty)
          (ps : List (Prod Ty Name)) (e : Expr) (sigma : TSubst),
          List.Mem dec d /\ dec.name = c /\
          List.Mem (Method.mk delta'' t0 m ps e) dec.methods /\
          mkSubst (dom dec.bounds ++ dom delta'') (ts ++ vars (dom delta'))
            = some sigma /\
          Ty.substList sigma (cod delta'') = cod delta' /\
          Ty.substList sigma (t0 :: ps.map Prod.fst) = t0' :: ps'.map Prod.fst)
    \/ (exists fs c0 ts0 es ms delta'' t0 ps e sigmaF sigma,
          Declared d (.loc c fs c0 ts0 es ms) /\
          List.Mem (Method.mk delta'' t0 m ps e) ms /\
          mkSubstF fs ts = some sigmaF /\
          mkSubst (dom delta'') (vars (dom delta')) = some sigma /\
          Ty.substList (sigmaF ++ sigma) (cod delta'') = cod delta' /\
          Ty.substList (sigmaF ++ sigma) (t0 :: ps.map Prod.fst)
            = t0' :: ps'.map Prod.fst)
  | .var _, _ => False
\end{lstlisting}
\end{leanside}
\mapping{The three cases are the three disjuncts. Comparing the lists $T_0,\ldots,T_n$ and $T'_0,\ldots,T'_n$ also requires the same number of parameters, as the paper does by writing $n$ in both methods. $[\pl{F},\Delta''=\ldots]$ is \texttt{sigmaF ++ sigma}, whose domains are disjoint by A7.}
\end{cmp}

\subsubsection{Methods, Classes and Programs}

\begin{cmp}{\textsc{method-ok}}
\begin{paperside}
\fit{$\inferrule{
  \kind{\Delta,\Delta'}{\mathrm{cod}(\Delta')}\\
  \kind{\Delta,\Delta'}{T_0,\ldots,T_n}\\\\
  \typing{\Delta,\Delta'}{\Gamma,T_1\gap x_1,\ldots,T_n\gap x_n}{e}{T_0}
}{
  \mok{\Delta}{\Gamma}{\method{\Delta'}{T_0}{m}{T_1\gap x_1,\ldots,T_n\gap x_n}{\ret{e}}}
}$}
\end{paperside}
\begin{leanside}
\begin{lstlisting}[style=leanside]
inductive MethodOk (d : Decls) (java : Bool) : Bounds -> Ctx -> Method -> Prop
  | mk {delta gamma delta' t0 m ps e} :
      (forall t, List.Mem t (cod delta') -> Kind d (delta ++ delta') t) ->
      (forall t, List.Mem t (t0 :: ps.map Prod.fst) -> Kind d (delta ++ delta') t) ->
      Typing d java (delta ++ delta') (gamma ++ ps.map fun p => (p.2, p.1)) e t0 ->
      MethodOk d java delta gamma (.mk delta' t0 m ps e)
\end{lstlisting}
\end{leanside}
\mapping{\texttt{MethodOk} is defined together with \texttt{Typing}, since \textsc{new-local} has a premise $\mok{\Delta'}{\ldots}{\pl{M}}$.}
\end{cmp}

\begin{cmp}{\textsc{class-ok}}
\begin{paperside}
\fit{$\inferrule{
  \kind{\Delta}{\mathrm{cod}(\Delta)}\\
  \kind{\Delta}{\gen{C'}{\pl{T}},T_1,\ldots,T_n}\\\\
  \mok{\Delta}{\gen{C}{\mathrm{dom}(\Delta)}\gap\this}{\pl{M}}\\\\
  \mathsf{overrideOk}(\Delta,\gen{C'}{\pl{T}},\pl{M})
}{
  \subclass{C}{\Delta}{C'}{\pl{T}}{\pl{M}}{T_1\gap f_1\semi\ldots T_n\gap f_n\semi}\gap\mathrm{Ok}
}$}
\end{paperside}
\begin{leanside}
\begin{lstlisting}[style=leanside]
def ClassOk (d : Decls) (java : Bool) (dec : Dec) : Prop :=
  (forall t, List.Mem t (cod dec.bounds) -> Kind d dec.bounds t) /\
  (forall t, List.Mem t (Ty.cls dec.sup dec.supArgs :: dec.fields.map Prod.fst) ->
    Kind d dec.bounds t) /\
  (forall mt, List.Mem mt dec.methods ->
    MethodOk d java dec.bounds [(thisName, .cls dec.name (vars (dom dec.bounds)))] mt) /\
  OverrideOk d dec.bounds (.cls dec.sup dec.supArgs) dec.methods
\end{lstlisting}
\end{leanside}
\mapping{A rule with one conclusion and no recursion is a definition in Lean: the four premises are the four conjuncts.}
\end{cmp}

\begin{cmp}{\textsc{prog-ok}}
\begin{paperside}
\fit{$\inferrule{
  \forall\mathsf{Dec}\in\decls\ldotp\mathsf{Dec}\gap\mathrm{Ok}\\\\
  \typing{\emptyset}{\emptyset}{e}{\_}
}{
  \vdash\decls\gap e\gap\mathrm{Ok}
}$}
\end{paperside}
\begin{leanside}
\begin{lstlisting}[style=leanside]
def ProgOk (d : Decls) (java : Bool) (e : Expr) : Prop :=
  (forall dec, List.Mem dec d -> ClassOk d java dec) /\
  exists t, Typing d java [] [] e t
\end{lstlisting}
\end{leanside}
\mapping{The $\_$ type is the existential \texttt{t}.}
\end{cmp}

\subsection{Statements and High Level Proofs}
\label{app:lean:statements}
\begin{lstlisting}[language=lean,basicstyle=\scriptsize\ttfamily]
def DeclsOk (d : Decls) : Prop := wfDecls d = true /\ forall dec, List.Mem dec d -> ClassOk d false dec
theorem stepWellFormed {d e e' t} (hd : DeclsOk d) (he : wfExpr e = true)
    (ht : Typing d false [] [] e t) (hs : Step d e e') : wfExpr e' = true
theorem progress {d e t} (hd : DeclsOk d) (he : wfExpr e = true)
    (ht : Typing d false [] [] e t) :
    e.isValue = true \/ exists e', Step d e e'
theorem preservation {d e e' t} (hd : DeclsOk d) (he : wfExpr e = true)
    (ht : Typing d false [] [] e t) (hs : Step d e e') : Typing d false [] [] e' t
theorem typeSoundness {d e0 e t} (hd : DeclsOk d) (he0 : wfExpr e0 = true)
    (ht : Typing d false [] [] e0 t) (hs : Steps d e0 e) :
    Typing d false [] [] e t /\ (e.isValue = true \/ exists e', Step d e e')
\end{lstlisting}
\texttt{progress}, \texttt{preservation} and \texttt{typeSoundness} are Progress, Preservation and Type Soundness of the section on the (un)soundness of method local classes.
\texttt{DeclsOk d} is the global \decls such that $\mathsf{Dec}\gap\mathrm{Ok}$ for every $\mathsf{Dec}\in\decls$, together with A1-A11 on \decls.
\texttt{stepWellFormed} has no counterpart in the paper (Section~\ref{app:lean:common}).
\texttt{typeSoundness} is proved by induction on \texttt{Steps}, from the other three theorems.
\texttt{progress} and \texttt{preservation} are proved from the low-level theorems with the same names, through these adequacy lemmas:
\begin{lstlisting}[language=lean,basicstyle=\scriptsize\ttfamily]
theorem isValue_tr {e : High.Expr} : (trExpr e).isValue = e.isValue
theorem wfExpr_tr {d : High.Decls} {e : High.Expr} (hd : High.wfDecls d = true)
    (he : High.wfExpr e = true) : Low.wfExpr (trExpr e) = true
theorem declsOk_tr {d : High.Decls} (hd : High.wfDecls d = true)
    (hok : forall dec, List.Mem dec d -> High.ClassOk d false dec) :
    Low.DeclsOk (trDecls d)
theorem typing_tr {d : High.Decls} {e : High.Expr} {t : High.Ty} (hd : High.wfDecls d = true) :
    High.Typing d false [] [] e t <-> Low.Typing (trDecls d) false [] [] (trExpr e) (trTy t)
theorem step_tr {d : High.Decls} {e e' : High.Expr} (hd : High.wfDecls d = true)
    (he : High.wfExpr e = true) (hs : High.Step d e e') :
    Low.Step (trDecls d) (trExpr e) (trExpr e')
theorem step_back {d : High.Decls} {e : High.Expr} {e'' : Low.Expr} (hd : High.wfDecls d = true)
    (he : High.wfExpr e = true) (hs : Low.Step (trDecls d) (trExpr e) e'') :
    exists e', High.Step d e e' /\ trExpr e' = e''
\end{lstlisting}
The adequacy lemmas, the translation and the low level can change together.
What must not change is the high level and the statements of \leanfile{Soundness.lean}.
While the two levels are identical, each adequacy lemma is an induction showing that the translation commutes with every function and preserves every judgement.
The planned proofs of the low-level \texttt{progress} and \texttt{preservation}, and of \texttt{stepWellFormed}, are:
\begin{enumerate}
\item \textbf{Basic lemmas.} \texttt{Ty.beq} is equality; \texttt{lookup} on a substitution built by \texttt{mkSubst}; by A1, \texttt{declaredKs}, \texttt{scope} and \texttt{generic} have one entry per class name.
\item \textbf{Renaming.} \texttt{Expr.freshen} renames bound names consistently, so typing, kinding and well-formedness are invariant under it; the result binds only names of \texttt{ns}.
\item \textbf{Weakening} of $\Gamma$ and $\Delta$, for typing, kinding and subtyping.
\item \textbf{Substitution of type variables.} If $\Delta,\Delta'$ types $e$, and ground types $\pl{G}$ satisfy the bounds $\Delta'$, then $e[\Delta'=\pl{G}]$ is typed under $\Delta$ with the substituted type.
Kinding and subtyping are preserved, but the checked cast is not: in rule \textsc{cast-var}, after $X$ becomes $G$, the source $G$ may not reach the target through \textsc{cast-up} or \textsc{cast-down}.
For the typing of rule \textsc{if-instanceof} the lemma uses rule \textsc{subs} to lift $e_0$ from $G$ to the substituted bound of $X$, and applies \textsc{cast-down} from there.
\item \textbf{Substitution of variables.} Values are closed and typed under $\emptyset;\emptyset$, so substituting values for variables of the same types preserves typing; \Q@this@ is not substituted into the methods of a local class.
\item \textbf{Inversion} of typing for each form of expression, through rule \textsc{subs}.
\item \textbf{Canonical forms.} A value of type $\gen{C}{\pl{T}}$ has a class that is a subclass of $C$.
\item \textbf{Copies of a local class.} By $\mathsf{consistent}$, the methods of a copy have the signatures of the declaration with the funnelled type variables substituted; so rules \textsc{call} and \textsc{sub-inherit-local}, which read the declaration, agree with rule \textsc{innerFound}, which runs the copy.
\item \textbf{Checked cast.} If $\emptyset\vdash\mathsf{typeOf}(v)<:\gen{C}{\_}$, $v$ has type $T_0$ and $\cast{\emptyset}{T_0}{\gen{C}{\pl{T}}}$, then $v$ has type $\gen{C}{\pl{T}}$.
The proof uses that the chain of superclasses is unique and that $\pl{X}\subseteq\ftv{\pl{T}'}$ fixes $\pl{T}$.
\item \textbf{stepWellFormed} by induction on \texttt{Step}: the fresh names keep A7, and the other requirements are local.
\item \textbf{progress} by induction on typing, with canonical forms; a call terminates its search at a class declaring the method, before \Q@Object@, since typing found the method in a superclass.
\item \textbf{preservation} by induction on \texttt{Step}, with the lemmas above.
\end{enumerate}

\subsection{Validation Examples}
\label{app:lean:examples}
\leanfile{Examples.lean} defines the class tables below and the examples, on the high level.
A Lean theorem marked \emph{proved} is accepted by Lean without \texttt{sorry}; a theorem marked \emph{stated} has \texttt{sorry}, and needs the inversion lemmas of Section~\ref{app:lean:statements}.
\begin{lstlisting}
class Object<> extends Object<>{}
class A<> extends Object<>{}
class B<> extends Object<>{ B<> onlyB(){ return new B<>(); } }
class Box<X extends Object<>> extends Object<>{ X val; X get(){ return this.val; } }
class C1<> extends Object<>{ A<> m(){ return new A<>(); } }
class C2<> extends C1<>{}                                         // d1: the classes above
class Foo<X extends Object<>> extends Object<>{ X x(){ return this.x(); } }
class Maker<> extends Object<>{ <Y extends Object<>> Foo<Y> mk(Y y){ return
  new MyFoo<Y funnel> extends Foo<Y>(){ Y x(){ return y; } }; } }  // d2: Object, A, B, Foo, Maker
class Broken<> extends Object<>{ <X extends Object<>> Foo<X> m(Object<> other, X x){ return
  if (other instanceof MyFoo<> w){ w }
  else { new MyFoo<> extends Foo<X>(){ X x(){ return x; } } }; } }  // d3: Object, A, B, Foo, Broken
class Leaker<> extends Object<>{
  <Y extends Object<>> MyFoo2<> mk2(Y y){ return new MyFoo2<> extends Foo<Y>(){ Y x(){ return y; } }; }
  <Y extends Object<>> Foo<Y> leak(MyFoo2<> z){ return z; } }     // d4: Object, A, B, Foo, Leaker
\end{lstlisting}
\begin{description}
\item[W1] \texttt{w1}, proved: d1 satisfies A1-A11.
\item[W2] \texttt{w2}, proved: a method with two parameters \Q@x@ violates A4.
\item[W3] \texttt{w3}, proved: an \Q@if@ binding the parameter \Q@x@ again violates A7.
\item[W4] \texttt{w4}, proved: an undeclared superclass violates A9 and A10.
\item[W5] \texttt{w5}, proved: d2 satisfies A1-A11.
\item[W6] \texttt{w6}, proved: d3 and the main expression of Example~1 satisfy A1-A11.
\item[W7] \texttt{w7}, proved: d4 satisfies A1-A11.
\item[R1] \texttt{r1}, proved: \Q@new Box<A<>>(new A<>()).val@ $\rightarrow$ \Q@new A<>()@ by \textsc{simpleFieldFound}.
\item[R2] \texttt{r2}, proved: \Q@new C2<>().m()@ $\rightarrow$ \Q@new C2<>().[C2<>]m()@ by \textsc{simpleAddSuper}.
\item[R3] \texttt{r3}, proved: \Q@new C2<>().[C2<>]m()@ $\rightarrow$ \Q@new C2<>().[C1<>]m()@ by \textsc{searchSuper}.
\item[R4] \texttt{r4}, proved: \Q@new C2<>().[C1<>]m()@ $\rightarrow$ \Q@new A<>()@ by \textsc{simpleFound}.
\item[R5] \texttt{r5}, proved: R1 inside the receiver of a call, by \textsc{ctx}.
\item[R6] \texttt{r6Then}, proved: \Q@if (new A<>() instanceof A<> x){ x } else { new B<>() }@ $\rightarrow$ \Q@new A<>()@ by \textsc{ifThen}.
\item[R6] \texttt{r6Else}, stated: the same with \Q@B<>@ as the test $\rightarrow$ \Q@new B<>()@ by \textsc{ifElse}.
\item[R7] \texttt{r7}, proved: \Q@new Maker<>().[Maker<>]<A<>>mk(new A<>())@ $\rightarrow$ \Q@new MyFoo<A<>> extends Foo<A<>>(){ A<> x(){ return new A<>(); } }@ by \textsc{simpleFound}, renaming \Q@Y@ and \Q@y@ into \Q@Y1@ and \Q@y1@.
\item[R8] \texttt{r8}, proved: the value of R7 \Q@.x()@ $\rightarrow$ \Q@new A<>()@ by \textsc{innerFound}.
\item[R9] \texttt{r9}, proved: the value of R7 \Q@.other()@ $\rightarrow$ \Q@.[Foo<A<>>]other()@ by \textsc{innerSuper}.
\item[R10] \texttt{r10}, proved: the value of R7 \Q@.f@ $\rightarrow$ \Q@new Foo<A<>>().f@ by \textsc{innerFieldSuper}.
\item[R11] \texttt{r11}, proved: R2, R3, R4 in sequence.
\item[T1] \texttt{t1}, proved: \Q@new Box<A<>>(new A<>()).val : A<>@.
\item[T2] \texttt{t2}, stated: \Q@new Box<A<>>(new B<>())@ is ill typed: \Q@B<>@ is not a subtype of \Q@A<>@.
\item[C1] \texttt{c1}, proved: $\cast{\QQ{X extends Object<>}}{\QQ{Foo<X>}}{\QQ{MyFoo<X>}}$ in d2.
\item[C2] \texttt{c2}, stated: not $\cast{\QQ{X extends Object<>}}{\QQ{Object<>}}{\QQ{MyFoo<X>}}$ in d2.
\item[M1] \texttt{m1}, stated: class \Q@Maker@ is Ok under funnelled typing.
\item[J1] \texttt{j1}, stated: Example~1 (d3) is well typed under Java typing.
\item[J2] \texttt{j2}, stated: Example~1 reduces to \Q@new A<>().onlyB()@, which is stuck.
\item[N1] \texttt{n1}, stated: Example~1 is ill typed under funnelled typing.
\item[L1] \texttt{l1}, stated: class \Q@Leaker@ is Ok under Java typing.
\item[L2] \texttt{l2}, stated: class \Q@Leaker@ is not Ok under funnelled typing.
\end{description}
\smallskip

\noindent In the paper, the verdict of each example follows from the rule named in the list, or, for W1-W7, from the requirements named in the list.
For R7, the paper renames with $\mathsf{fresh}$ the type variable \Q@Y@ and the parameter \Q@y@ of \Q@mk@, and then substitutes \Q@A<>@ and \Q@new A<>()@; the method \Q@x@ of the local class binds no names.
For T2, rule \textsc{new-top} needs \Q@new B<>() : A<>@, and no rule gives it.
For C2, rule \textsc{cast-down} needs $\QQ{X}\in\ftv{\emptyset}$, since \Q@MyFoo<Y>@ has the supertype \Q@Object<>@ with no type arguments.
For J1, J2 and N1, Section~\ref{sec:java} gives the derivation; N1 fails at $\kind{\emptyset}{\QQ{Foo<X>}}$ in rule \textsc{new-local}, since \Q@MyFoo@ funnels nothing.
For L1 and L2, Section~\ref{sec:java} discusses the method \Q@leak@; L2 fails at the local class in \Q@mk2@, as N1 does.
